% answers.tex —— 1985 年·上海卷·文史类 · 参考答案与详细解析片段
% 由 tools/generate_projects.py 依来源总账自 1985/1985_数学真题_上海文_试题.md + 1985/1985_数学真题_上海文_解析正文.md + 1985/1985_数学真题_上海文_答案要点.md 生成，勿手改（改源 MD 后重跑）


\begin{gktip}
固定源为两页题面，不含答案页。下列逐题推导均为本库独立复算；第 5、8 题的“如图”未见独立源图，解析不引入不可核验的图形条件。
\end{gktip}

\begin{answer}
\gksolhead{1}{}

\textbf{题目}：

（本题满分 30 分）

本题共有 10 个小题，每一个小题满分 3 分，只要求直接填写结果。

\Needspace{5\baselineskip}
（1） 不等式

\nopagebreak[4]
\[
|x+2|<1
\]

的解集是　　　　。

\Needspace{5\baselineskip}
（2） 函数

\nopagebreak[4]
\[
y=\sqrt[8]{x}-1
\]

的反函数是　　　　。

\Needspace{5\baselineskip}
（3） 函数

\nopagebreak[4]
\[
y=\operatorname{tg}2x
\]

的最小正周期是　　　　。

（4） 设等差数列 $\{a_n\}$ 中，$a_1=2$，$a_2=5$，则
$a_{10}=$　　　　。

\Needspace{5\baselineskip}
（5） 方程

\nopagebreak[4]
\[
9^x=3^{x+1}
\]

的解是　　　　。

\Needspace{5\baselineskip}
（6） 用不等号“$>$”或“$<$”连接：

\nopagebreak[4]
\[
\left(\frac12\right)^{\frac17}\ \underline{\hspace{2em}}\ 1.
\]

（7） 如果一个圆的圆心在点 $(2,4)$，并且经过点 $(0,3)$，则这个圆的方程是　　　　。

\Needspace{5\baselineskip}
（8） 直线

\nopagebreak[4]
\[
y=-x+3
\]

的倾斜角的大小是　　　　。

\Needspace{5\baselineskip}
（9） 设 $\lg2=a$，则

\nopagebreak[4]
\[
\log_2 25=
\]

　　　　。

（10） 已知 $O$ 为直角坐标系的原点，$A$ 点的坐标为 $(1,3)$，$B$ 点的坐标为 $(0,3)$，若把 $\triangle OAB$ 绕 $y$ 轴旋转一周，则所得旋转体的体积是　　　　。

\Needspace{4\baselineskip}
\textbf{解析}

\gksubpart{（1）}

\Needspace{5\baselineskip}
由 $|x+2|<1$ 得

\nopagebreak[4]
\[
-1<x+2<1,
\]

\Needspace{5\baselineskip}
所以

\nopagebreak[4]
\[
-3<x<-1.
\]

\gksubpart{（2）}

\Needspace{5\baselineskip}
由

\nopagebreak[4]
\[
y=\sqrt[8]{x}-1
\]

\Needspace{5\baselineskip}
得 $y+1=\sqrt[8]{x}$，即 $x=(y+1)^8$。交换 $x\text{、}y$，反函数为

\nopagebreak[4]
\[
y=(x+1)^8,\qquad x\ge-1.
\]

\gksubpart{（3）}

\Needspace{5\baselineskip}
$\operatorname{tg}u$ 的最小正周期为 $\pi$，令 $u=2x$，故

\nopagebreak[4]
\[
T=\frac\pi2.
\]

\gksubpart{（4）}

\Needspace{5\baselineskip}
等差数列公差 $d=a_2-a_1=3$，所以

\nopagebreak[4]
\[
a_{10}=a_1+9d=2+27=29.
\]

\gksubpart{（5）}

\nopagebreak[4]
\[
9^x=3^{2x},\qquad 3^{2x}=3^{x+1},
\]

故 $2x=x+1$，解得 $x=1$。

\gksubpart{（6）}

\Needspace{5\baselineskip}
因为 $0<\frac12<1$，其正分数次幂仍小于 1，所以

\nopagebreak[4]
\[
\left(\frac12\right)^{1/7}<1.
\]

\gksubpart{（7）}

\Needspace{5\baselineskip}
半径平方为

\nopagebreak[4]
\[
r^2=(2-0)^2+(4-3)^2=5,
\]

\Needspace{5\baselineskip}
圆方程为

\nopagebreak[4]
\[
(x-2)^2+(y-4)^2=5.
\]

\gksubpart{（8）}

\Needspace{5\baselineskip}
直线斜率为 $-1$。倾斜角 $\alpha\in[0,\pi)$ 满足 $\tan\alpha=-1$，故

\nopagebreak[4]
\[
\alpha=\frac{3\pi}{4}.
\]

\gksubpart{（9）}

\nopagebreak[4]
\[
\log_2 25=2\log_2 5
=2\frac{\lg5}{\lg2}
=\frac{2(1-a)}a.
\]

\gksubpart{（10）}

\Needspace{5\baselineskip}
底面三角形 $OAB$ 绕 $y$ 轴旋转形成底面半径为 $1$、高为 $3$ 的圆锥，因此

\nopagebreak[4]
\[
V=\frac13\pi\cdot1^2\cdot3=\pi.
\]

\gkanswer{答案（1）：$(-3,-1)$}

\gkanswer{答案（2）：$y=(x+1)^8\ (x\ge-1)$}

\gkanswer{答案（3）：$\dfrac\pi2$}

\gkanswer{答案（4）：$29$}

\gkanswer{答案（5）：$x=1$}

\gkanswer{答案（6）：$\left(\dfrac12\right)^{1/7}<1$}

\gkanswer{答案（7）：$(x-2)^2+(y-4)^2=5$}

\gkanswer{答案（8）：$\dfrac{3\pi}{4}$}

\gkanswer{答案（9）：$\dfrac{2(1-a)}a$}

\gkanswer{答案（10）：$\pi$}

\end{answer}
\begin{answer}
\gksolhead{2}{}

\textbf{题目}：

（本题满分 15 分）

本题共有 5 个小题，每一个小题都给出代号为 A、B、C、D 的四个结论，其中只有一个结论是正确的。把正确结论的代号写在题后的圆括号内，选对得 3 分，不选、选错或者选出的代号超过一个（不论是否都写在圆括号内），一律得 0 分。

\Needspace{5\baselineskip}
（1） 函数

\nopagebreak[4]
\[
y=-\sin2x+3
\]

的最大值是\mbox{（　　）}

（A）2　　（B）3　　（C）4　　（D）5

（2） 设复数 $z_1=1+i$，$z_2=1-3i$，则复数
$z=z_1z_2$ 在复平面内所表示的点位于\mbox{（　　）}

（A）第一象限　　（B）第二象限　　（C）第三象限　　（D）第四象限

\Needspace{5\baselineskip}
（3） 已知函数

\nopagebreak[4]
\[
f(x)=\lg x^2
\]

\Needspace{5\baselineskip}
的定义域为 $F$，函数

\nopagebreak[4]
\[
g(x)=2\lg x
\]

的定义域为 $G$，那么\mbox{（　　）}

（A）$F\cap G=\varnothing$　　（B）$F=G$

（C）$F\subset G$　　（D）$G\subset F$

（4） 若一个四棱锥的底面是凸四边形，它的顶点到底面各边的距离都相等，则这个棱锥的底面四边形\mbox{（　　）}

（A）必为正方形　　（B）必有内切圆

（C）必有外接圆　　（D）必既有内切圆又有外接圆

（5） $a>0$ 且 $b>0$，是 $ab>0$ 的\mbox{（　　）}

（A）充分但不必要的条件　　（B）必要但不充分的条件

（C）充分而且必要的条件　　（D）既不充分又不必要的条件

\Needspace{4\baselineskip}
\textbf{解析}

\gksubpart{（1）}

因为 $-\sin2x\le1$，所以 $y=-\sin2x+3\le4$；当
$\sin2x=-1$ 时取到 4。故选（C）。

\gksubpart{（2）}

\nopagebreak[4]
\[
z=(1+i)(1-3i)=4-2i,
\]

实部正、虚部负，点在第四象限，选（D）。

\gksubpart{（3）}

\nopagebreak[4]
\[
F=\mathbb R\setminus\{0\},\qquad G=(0,+\infty).
\]

因此 $G\subset F$，选（D）。

\gksubpart{（4）}

设棱锥顶点为 $V$，$V$ 在底面的射影为 $O$。点 $V$ 到底面各边的距离相等，等价于 $O$ 到四条底边所在直线的距离相等；由于底面是凸四边形，$O$ 在其内部，从而底面四边形有内切圆。选（B）。

\gksubpart{（5）}

$a>0$ 且 $b>0$ 能推出 $ab>0$，但 $ab>0$ 也可能由 $a<0\text{、}b<0$ 得到。因此这是充分但不必要的条件，选（A）。

\gkanswer{答案（1）：（C）}

\gkanswer{答案（2）：（D）}

\gkanswer{答案（3）：（D）}

\gkanswer{答案（4）：（B）}

\gkanswer{答案（5）：（A）}

\end{answer}
\begin{answer}
\gksolhead{3}{}

\textbf{题目}：

（本题满分 9 分）

本题共有 2 个小题，第（1）小题满分 4 分，第（2）小题满分 5 分。

在平面直角坐标系内，画出下列方程的曲线（只要求画出图形）：

\Needspace{5\baselineskip}
（1）　

\nopagebreak[4]
\[
\frac{x}{3}-\frac{y}{2}=-1;
\]

\Needspace{5\baselineskip}
（2）　

\nopagebreak[4]
\[
\frac{x^2}{9}-\frac{y^2}{4}=-1.
\]
\nopagebreak[4]
\begin{center}% 1985 年上海卷文史类第 3 题：直线与竖轴双曲线
\begin{tikzpicture}[scale=0.72,>=stealth]
  \begin{scope}[shift={(-3.0,0)}]
    \draw[->] (-2.1,0)--(2.25,0) node[right] {$x$};
    \draw[->] (0,-1.0)--(0,2.8) node[above] {$y$};
    \draw[thick] (-2.0,0.67)--(1.2,2.8);
    \fill (0,2) circle (1.1pt);
    \node[below left] at (0,0) {$O$};
    \node[below] at (0,-1.25) {（1）};
  \end{scope}
  \begin{scope}[shift={(2.0,0)}]
    \draw[->] (-1.85,0)--(1.9,0) node[right] {$x$};
    \draw[->] (0,-2.4)--(0,2.4) node[above] {$y$};
    \draw[thick,domain=-1.65:1.65,samples=100]
      plot (\x,{2*sqrt(1+\x*\x/9)});
    \draw[thick,domain=-1.65:1.65,samples=100]
      plot (\x,{-2*sqrt(1+\x*\x/9)});
    \fill (0,2) circle (1.1pt);
    \fill (0,-2) circle (1.1pt);
    \node[below left] at (0,0) {$O$};
    \node[below] at (0,-2.75) {（2）};
  \end{scope}
\end{tikzpicture}
\end{center}
\nopagebreak[4]
\Needspace{4\baselineskip}
\textbf{解析}

\gksubpart{（1）}

\Needspace{5\baselineskip}
由

\nopagebreak[4]
\[
\frac{x}{3}-\frac{y}{2}=-1
\]

\Needspace{5\baselineskip}
得

\nopagebreak[4]
\[
y=\frac23x+2.
\]

这是斜率为 $\frac23$、纵截距为 2 的直线。

\gksubpart{（2）}

\Needspace{5\baselineskip}
移项得

\nopagebreak[4]
\[
\frac{y^2}{4}-\frac{x^2}{9}=1,
\]

这是中心在原点、焦点在 $y$ 轴方向的等轴双曲线，顶点为
$(0,\pm2)$，渐近线为 $y=\pm\frac23x$。

\gkanswer{答案（1）：直线 $y=\dfrac23x+2$。}

\gkanswer{答案（2）：双曲线 $\dfrac{y^2}{4}-\dfrac{x^2}{9}=1$。}

\end{answer}
\begin{answer}
\gksolhead{4}{}

\textbf{题目}：

（本题满分 20 分）

本题共有 3 个小题，第（1）小题满分 5 分，第（2）小题满分 7 分，第（3）小题满分 8 分。

\Needspace{5\baselineskip}
（1） 求

\nopagebreak[4]
\[
(2+x^3)^8
\]

的展开式中 $x^{10}$ 的系数。

（2） 从六个数字 $1,2,3,5,7,9$ 中任取四个不同的数字，有多少种取法？由这六个数字可以组成多少个没有重复数字的四位偶数？

（3） 在一列球中，第 1 个球的半径为 1，第 2 个球的直径是第 1 个球的半径，以后每一个球的直径都是前一个球的半径，这样无限继续下去，求所有这些球的表面积的和。

\Needspace{4\baselineskip}
\textbf{解析}

\gksubpart{（1）}

\Needspace{5\baselineskip}
二项式通项为

\nopagebreak[4]
\[
\binom8k2^{8-k}x^{3k}.
\]

其中 $x$ 的指数只能是 3 的倍数，不能出现 $x^{10}$ 项，所以系数为 0。

\gksubpart{（2）}

\Needspace{5\baselineskip}
从 6 个数字中任取 4 个，取法数为

\nopagebreak[4]
\[
\binom64=15.
\]

\Needspace{5\baselineskip}
要组成四位偶数，末位必须是唯一的偶数 2。其余三位从另外 5 个数字中选出并排列，故共有

\nopagebreak[4]
\[
5\cdot4\cdot3=60
\]

个。

\gksubpart{（3）}

\Needspace{5\baselineskip}
第 $n$ 个球的半径为

\nopagebreak[4]
\[
r_n=\left(\frac12\right)^{n-1}.
\]

\Needspace{5\baselineskip}
因此所有球表面积之和为

\nopagebreak[4]
\[
\begin{aligned}
\sum_{n=1}^{\infty}4\pi r_n^2
&=4\pi\sum_{n=1}^{\infty}\left(\frac14\right)^{n-1}\\
&=\frac{4\pi}{1-\frac14}
=\frac{16\pi}{3}.
\end{aligned}
\]

\gkanswer{答案（1）：$0$}

\gkanswer{答案（2）：15 种取法，60 个四位偶数}

\gkanswer{答案（3）：$\dfrac{16\pi}{3}$}

\end{answer}
\begin{answer}
\gksolhead{5}{}

\textbf{题目}：

（本题满分 10 分）

如图，直线 $a\parallel$ 直线 $b$，点 $P$ 在 $a\text{、}b$ 所确定的平面外，$PA\perp a$，且交 $a$ 于 $A$ 点；$AB\perp b$，且交 $b$ 于 $B$ 点。

求证：$PB\perp b$。

\begin{gktip}
固定源页面保留“如图”题干，但未见可独立辨认的配图；证明仅依据可见直线关系转录。
\end{gktip}

\Needspace{4\baselineskip}
\textbf{解析}

因 $a\parallel b$，且 $PA\perp a$，故 $PA$ 的方向垂直于
\Needspace{5\baselineskip}
$b$ 的方向；又 $AB\perp b$。直线 $PA\text{、}AB$ 相交于 $A$，且都垂直于 $b$ 的方向，所以平面 $PAB$ 垂直于直线 $b$。由于 $PB$ 在平面 $PAB$ 内，并交 $b$ 于 $B$，故

\nopagebreak[4]
\[
PB\perp b.
\]

\gkanswer{答案：$PB\perp b$。}

\end{answer}
\begin{answer}
\gksolhead{6}{}

\textbf{题目}：

（本题满分 12 分）

本题共有 2 个小题，第（1）小题满分 4 分，第（2）小题满分 8 分。

\Needspace{5\baselineskip}
（1） 证明万能公式：

\nopagebreak[4]
\[
\sin\alpha=\frac{2\operatorname{tg}\frac\alpha2}
{1+\operatorname{tg}^2\frac\alpha2}.
\]

\Needspace{5\baselineskip}
（2） 已知

\nopagebreak[4]
\[
\cos\alpha=-\frac13,\qquad
\pi<\alpha<\frac{3\pi}{2},
\]

\Needspace{5\baselineskip}
求

\nopagebreak[4]
\[
\cos\left(\frac\alpha2+\frac{5\pi}{6}\right).
\]

\Needspace{4\baselineskip}
\textbf{解析}

\gksubpart{（1）}

\Needspace{5\baselineskip}
由二倍角公式

\nopagebreak[4]
\[
\sin\alpha=2\sin\frac\alpha2\cos\frac\alpha2
\]

\Needspace{5\baselineskip}
以及

\nopagebreak[4]
\[
1+\operatorname{tg}^2\frac\alpha2
=\frac1{\cos^2\frac\alpha2},
\qquad
\operatorname{tg}\frac\alpha2
=\frac{\sin\frac\alpha2}{\cos\frac\alpha2},
\]

\Needspace{5\baselineskip}
可得

\nopagebreak[4]
\[
\frac{2\operatorname{tg}\frac\alpha2}
{1+\operatorname{tg}^2\frac\alpha2}
=\frac{2\frac{\sin\frac\alpha2}{\cos\frac\alpha2}}
{\frac1{\cos^2\frac\alpha2}}
=2\sin\frac\alpha2\cos\frac\alpha2
=\sin\alpha.
\]

\gksubpart{（2）}

\Needspace{5\baselineskip}
由 $\pi<\alpha<\frac{3\pi}{2}$，知 $\frac\alpha2$ 位于第二象限。因此

\nopagebreak[4]
\[
\cos\frac\alpha2=-\sqrt{\frac{1+\cos\alpha}{2}}
=-\frac1{\sqrt3},
\qquad
\sin\frac\alpha2=\sqrt{\frac{1-\cos\alpha}{2}}
=\sqrt{\frac23}.
\]

\Needspace{5\baselineskip}
于是

\nopagebreak[4]
\[
\begin{aligned}
\cos\left(\frac\alpha2+\frac{5\pi}{6}\right)
&=\cos\frac\alpha2\cos\frac{5\pi}{6}
-\sin\frac\alpha2\sin\frac{5\pi}{6}\\
&=\frac12-\frac12\sqrt{\frac23}\\
&=\frac{\sqrt3-\sqrt2}{2\sqrt3}.
\end{aligned}
\]

\gkanswer{答案（1）：见解析正文。}

\gkanswer{答案（2）：$\dfrac{\sqrt3-\sqrt2}{2\sqrt3}$}

\end{answer}
\begin{answer}
\gksolhead{7}{}

\textbf{题目}：

（本题满分 10 分）

设 $\triangle ABC$ 的两个内角 $A\text{、}B$ 所对的边分别为 $a\text{、}b$。若复数 $z_1$ 的模为 $a$，幅角为 $B$；复数 $z_2$ 的模为 $b$，幅角为 $-A$。求证：$z_1+z_2$ 为实数。

\Needspace{4\baselineskip}
\textbf{解析}

\Needspace{5\baselineskip}
将两个复数写成三角形式：

\nopagebreak[4]
\[
z_1=a(\cos B+i\sin B),\qquad
z_2=b(\cos A-i\sin A).
\]

\Needspace{5\baselineskip}
由正弦定理，

\nopagebreak[4]
\[
\frac a{\sin A}=\frac b{\sin B},
\]

\Needspace{5\baselineskip}
所以 $a\sin B=b\sin A$。因此 $z_1+z_2$ 的虚部为

\nopagebreak[4]
\[
a\sin B-b\sin A=0,
\]

故 $z_1+z_2$ 为实数。

\gkanswer{答案：见解析正文。}

\end{answer}
\begin{answer}
\gksolhead{8}{}

\textbf{题目}：

（本题满分 14 分）

\Needspace{5\baselineskip}
如图，由抛物线

\nopagebreak[4]
\[
y=x^2+2
\]

\Needspace{5\baselineskip}
上的点 $M(x_0,y_0)$ 向直线

\nopagebreak[4]
\[
y=\frac12x
\]

\Needspace{5\baselineskip}
作垂线，垂足为 $N$，延长 $MN$ 至 $P$ 点，使点 $N$ 分线段 $MP$ 成定比

\nopagebreak[4]
\[
\lambda=\frac{MN}{NP}=4.
\]

（1） 用 $x_0,y_0$ 表示点 $N$ 的坐标 $(x_1,y_1)$；

（2） 用 $x_0,y_0$ 和 $x_1,y_1$ 表示点 $P$ 的坐标 $(x,y)$；

（3） 求：当点 $M$ 沿抛物线移动时，点 $P$ 的轨迹方程，并说明轨迹是什么图形。

\begin{gktip}
固定源页面保留“如图”题干，但未见可独立辨认的配图；本题坐标推导仅依据可见文字转录。
\end{gktip}

\Needspace{4\baselineskip}
\textbf{解析}

\gksubpart{（1）}

直线 $y=\frac12x$ 的方向向量可取 $(2,1)$。点
\Needspace{5\baselineskip}
$M(x_0,y_0)$ 在该方向上的投影为

\nopagebreak[4]
\[
N=\frac{2x_0+y_0}{5}(2,1),
\]

\Needspace{5\baselineskip}
所以

\nopagebreak[4]
\[
x_1=\frac{4x_0+2y_0}{5},
\qquad
y_1=\frac{2x_0+y_0}{5}.
\]

\gksubpart{（2）}

由于 $M\text{、}N\text{、}P$ 共线，且 $MN/NP=4$，点 $N$ 到 $P$ 的方向与
\Needspace{5\baselineskip}
$M$ 到 $N$ 的方向相同，并有

\nopagebreak[4]
\[
\overrightarrow{NP}=\frac14\overrightarrow{MN}.
\]

\Needspace{5\baselineskip}
因此

\nopagebreak[4]
\[
P=N+\frac14(N-M)=\frac{5N-M}{4},
\]

\Needspace{5\baselineskip}
从而

\nopagebreak[4]
\[
x=\frac{5x_1-x_0}{4}=\frac{3x_0+2y_0}{4},
\qquad
y=\frac{5y_1-y_0}{4}=\frac{x_0}{2}.
\]

\gksubpart{（3）}

由 $y=x_0/2$ 得 $x_0=2y$。又因 $M$ 在
\Needspace{5\baselineskip}
$y=x^2+2$ 上，

\nopagebreak[4]
\[
y_0=x_0^2+2=4y^2+2.
\]

\Needspace{5\baselineskip}
代入 $x=(3x_0+2y_0)/4$：

\nopagebreak[4]
\[
x=\frac{6y+8y^2+4}{4}
=2y^2+\frac32y+1.
\]

\Needspace{5\baselineskip}
所以轨迹普通方程为

\nopagebreak[4]
\[
x^2-y^2+2x-2y+8=0,
\]

\Needspace{5\baselineskip}
即

\nopagebreak[4]
\[
(y+1)^2-(x+1)^2=8.
\]

这是一条以 $(-1,-1)$ 为中心的等轴双曲线。

\gkanswer{答案（1）：$x_1=\dfrac{4x_0+2y_0}{5}$，$y_1=\dfrac{2x_0+y_0}{5}$}

\gkanswer{答案（2）：$x=\dfrac{3x_0+2y_0}{4}$，$y=\dfrac{x_0}{2}$}

\gkanswer{答案（3）：$x^2-y^2+2x-2y+8=0$，即 $(y+1)^2-(x+1)^2=8$，是一条等轴双曲线。}

\end{answer}